% !PRINTED 597
\weberpage{597}
According to this, the highest power of the prime ideal $\mathfrak{p}$ that divides the partial fundamental ideal $\mathfrak{G}_R$ can be given. By \S.\,159, namely, this ideal is defined by the product
\[ f'(\tau) = (\tau - \tau_1)(\tau - \tau_2)\ldots(\tau - \tau_{n-1}), \]
in which $\tau_1, \tau_2, \ldots, \tau_{n-1}$ are those among the functionals conjugate to $\tau$ that are different from $\tau$. One of the factors of this product, $\tau - \tau\mkern1mu|\mkern1mu\varphi$, is then and only then divisible by $\mathfrak{p}$, when $\varphi$ belongs to $X$. Now since the identical substitution is excluded, it follows that $\mathfrak{G}$ contains the factor $\mathfrak{p}$ at least $g-1$ times.

But $\tau - \tau\mkern1mu|\mkern1mu\chi$ is then and only then divisible by $\mathfrak{p}^2$, when $\chi$ belongs to $X^{(1)}$, and consequently $\mathfrak{G}$ is divisible by $\mathfrak{p}^{g-1+g_1-1}$. We can continue in this way and find the exact power of $\mathfrak{p}$ that is contained in $\mathfrak{G}$. It has the exponent
\[ (g-1) + (g_1-1) + (g_2-1) + \cdots \]
Each of the numbers $g, g_1, g_2, \ldots$ is a divisor of the one preceding, and $g_1, g_2, \ldots$ are powers of $p$, and the sequence continues until one of the numbers has become $g = 1$. The composition of the fundamental ideal $\mathfrak{G}$ is thus completely determined by the decomposition of the group $X$.

The numbers $g, g_1, g_2, \ldots$ are moreover the same for all the conjugate prime ideals $\mathfrak{p}_1, \mathfrak{p}_2, \ldots, \mathfrak{p}_e$, and it therefore follows that
\[ \mathfrak{G} = \Pi(\mathfrak{p}_1\mathfrak{p}_2\ldots\mathfrak{p}_e)^{g-1+g_1-1+g_2-2\cdots}, \] % [sic] exponent pattern breaks from g_2-1 to g_2-2 as printed
where the product sign refers to the prime factors of all the primes $p$ of the field $\Omega$. In this, of course, only those primes $p$, present in finite number, for which $g > 1$ need be taken into account.

The partial norm of the ideal $\mathfrak{G}$
\[ \mathfrak{D} = \mathfrak{N}_R(\mathfrak{G}) = \Pi\, \mathfrak{P}^{ef(g-1+g_1-1+g_2-1+\cdots)} \]
is called the partial discriminant of the field $\Omega$ with respect to the field $R$. It is an ideal lying in the field $R$, which, in the case that $R$ is the field of rational numbers, passes over, apart from a unit factor, into the fundamental number of the field $\Omega$ (\S.\,158).

\begin{center}
\S.\,163.

\gesperrt{The Ideals in the Divisors of the Field $\Omega$.}
\end{center}

If the group $\Phi$ of the field $\Omega$ has a divisor $\Phi'$, then to $\Phi'$ there belongs a field $\Omega'$, which is a divisor of $\Omega$ and
%CONT>>
